% !TeX root = ../main.tex
% Fonte letta integralmente: Lezione 1.pdf, pp. 2--6.
% Le spiegazioni, la tabella di verità e gli esempi sviluppano i passaggi
% sintetici della fonte senza attribuire al docente formulazioni aggiunte.
\chapter{Logica e metodi di dimostrazione}
\label{cap:logica}
\fonti{Lezione 1}

\section{Proposizioni e connettivi logici}
% Fonte: Lezione 1, pp.~2--3

Una \emph{proposizione} è un'affermazione alla quale si può attribuire un
valore di verità: vero oppure falso. Per esempio, «$2+3=5$» è una
proposizione vera, mentre «$2+3=6$» è una proposizione falsa. Una frase
contenente una variabile libera, come «$x>0$», richiede invece di sapere
quale valore si assegna a $x$ prima di poter stabilire se sia vera o falsa.

Per costruire proposizioni complesse a partire da proposizioni $P$ e $Q$
si usano i \emph{connettivi logici}:
\begin{description}[style=nextline]
  \item[Negazione: $\neg P$.]
    È vera esattamente quando $P$ è falsa; si legge «non $P$».
  \item[Congiunzione: $P\land Q$.]
    È vera quando sono vere sia $P$ sia $Q$; si legge «$P$ e $Q$».
  \item[Disgiunzione: $P\lor Q$.]
    È vera quando è vera almeno una delle due proposizioni. Si tratta
    della disgiunzione inclusiva: è vera anche se sono vere entrambe.
  \item[Implicazione: $P\Rightarrow Q$.]
    È falsa soltanto quando $P$ è vera e $Q$ è falsa. $P$ è l'ipotesi
    e $Q$ la tesi; si legge «se $P$, allora $Q$».
  \item[Doppia implicazione: $P\Leftrightarrow Q$.]
    È vera quando $P$ e $Q$ hanno lo stesso valore di verità: entrambe
    vere oppure entrambe false. Si legge «$P$ se e solo se $Q$».
\end{description}

La tabella seguente riassume tutti i casi; $V$ significa vero e $F$ falso.
\begin{center}
\renewcommand{\arraystretch}{1.2}
\begin{tabular}{ccccccc}
\toprule
$P$ & $Q$ & $\neg P$ & $P\land Q$ & $P\lor Q$ & $P\Rightarrow Q$ & $P\Leftrightarrow Q$\\
\midrule
$V$ & $V$ & $F$ & $V$ & $V$ & $V$ & $V$\\
$V$ & $F$ & $F$ & $F$ & $V$ & $F$ & $F$\\
$F$ & $V$ & $V$ & $F$ & $V$ & $V$ & $V$\\
$F$ & $F$ & $V$ & $F$ & $F$ & $V$ & $V$\\
\bottomrule
\end{tabular}
\end{center}

\section{Implicazioni e metodi di dimostrazione}
% Fonte: Lezione 1, p.~3

L'implicazione $P\Rightarrow Q$ esprime che la verità di $P$ basta a
garantire la verità di $Q$. Per questo $P$ è detta \emph{condizione
sufficiente} per $Q$ e $Q$ \emph{condizione necessaria} per $P$.
L'implicazione non afferma che $P$ sia vera: esclude soltanto il caso
in cui l'ipotesi sia vera e la tesi falsa. In particolare,
\[
  P\Rightarrow Q\quad\Longleftrightarrow\quad \neg P\lor Q.
\]
Per dimostrare una doppia implicazione si devono verificare entrambi
i versi, $P\Rightarrow Q$ e $Q\Rightarrow P$.

\begin{proposizione}[Contrapposizione]
Un'implicazione è equivalente alla sua contronominale:
\[
  (P\Rightarrow Q)\quad\Longleftrightarrow\quad
  (\neg Q\Rightarrow\neg P).
\]
\end{proposizione}
\begin{proof}
La prima implicazione è falsa esattamente nel caso $P$ vera e $Q$ falsa.
La seconda è falsa esattamente quando $\neg Q$ è vera e $\neg P$ è falsa,
cioè nello stesso caso. Le due implicazioni hanno quindi sempre lo
stesso valore di verità.
\end{proof}

Questa equivalenza fornisce un metodo di dimostrazione: se è più semplice
partire dalla negazione della tesi, si dimostra $\neg Q\Rightarrow\neg P$
e si conclude $P\Rightarrow Q$.

Il \emph{ragionamento per assurdo} consiste nell'assumere la negazione di ciò che si
vuole dimostrare e nel ricavarne una contraddizione. Per provare
$P\Rightarrow Q$ si può dunque assumere che $P$ sia vera e $Q$ falsa e
mostrare che queste due assunzioni non possono valere contemporaneamente.
I due metodi sono strettamente collegati, ma conviene distinguerne
il modo di procedere.

\section{Predicati e quantificatori}
% Fonte: Lezione 1, pp.~4--5

Un \emph{predicato} $P(x)$ è un'affermazione che dipende da una variabile
$x$. Il suo significato va accompagnato dall'insieme nel quale $x$ può
variare: per esempio, $P(x)$ può essere la proprietà «$x^2=1$» con
$x\in\R$. I quantificatori permettono di trasformare un predicato
in una proposizione.

\begin{definizione}[Quantificatore universale]
La proposizione
\[
  \forall x\in U,\quad P(x)
\]
è vera se $P(a)$ è vera per ogni valore $a$ appartenente all'insieme
ambiente $U$. Il simbolo $\forall$ si legge «per ogni» o «qualunque sia».
\end{definizione}

\begin{definizione}[Quantificatore esistenziale]
La proposizione
\[
  \exists x\in U:\ P(x)
\]
è vera se esiste almeno un elemento $a\in U$ per il quale $P(a)$ è vera.
Il simbolo $\exists$ si legge «esiste». Non si esclude che gli elementi
con questa proprietà siano più di uno.
\end{definizione}

\begin{definizione}[Esistenza e unicità]
La proposizione
\[
  \exists!\,x\in U:\ P(x)
\]
è vera quando esiste esattamente un elemento di $U$ che soddisfa $P$.
Per dimostrarla occorre provare sia l'esistenza sia l'unicità: dopo
aver trovato un elemento $a$ con la proprietà richiesta, si mostra che
ogni altro elemento che la soddisfa coincide con $a$.
\end{definizione}

Per esempio, $\exists x\in\R:x^2=1$ è vera, ma
$\exists!\,x\in\R:x^2=1$ è falsa, perché sia $1$ sia $-1$ risolvono
l'equazione. Se imponiamo anche $x>0$, la soluzione diventa unica.

\section{Negare una proposizione con quantificatori}
% Fonte: Lezione 1, p.~6

Negare un'affermazione universale significa trovare almeno un caso
in cui la proprietà non vale; negare un'affermazione esistenziale
significa affermare che la proprietà non vale per nessun elemento.
Formalmente,
\begin{align*}
  \neg\bigl(\forall x\in U,\ P(x)\bigr)
    &\quad\Longleftrightarrow\quad \exists x\in U:\ \neg P(x),\\
  \neg\bigl(\exists x\in U:\ P(x)\bigr)
    &\quad\Longleftrightarrow\quad \forall x\in U,\ \neg P(x).
\end{align*}
Quando la negazione attraversa un quantificatore, dunque, $\forall$
si trasforma in $\exists$ e $\exists$ si trasforma in $\forall$;
anche la proprietà all'interno deve essere negata.

Per esempio, la negazione di «tutti gli elementi di $U$ sono positivi»
è «esiste un elemento di $U$ che è minore o uguale a zero».
Non è necessario che siano tutti negativi: basta un controesempio.
Se sono presenti più quantificatori, si applica la regola a ciascuno,
mantenendo il loro ordine. Per esempio,
\[
  \neg\bigl(\forall x\in\R\ \exists y\in\R:\ y>x\bigr)
  \quad\Longleftrightarrow\quad
  \exists x\in\R\ \forall y\in\R:\ y\le x.
\]
